%% kapa: exam=1 round=28 year=2017 period=2 form=A subject=accounting subject_ko=회계학 seq=34 number=74
%% source: Q-Net 감정평가사 기출문제 (cst00309) artlSeq=5199603 file=extracted/1008/5199603/2017년 제28회 감정평가사 1차시험 문제지/2교시 A형.pdf
%% conversion: kordoc 4.24.51 -> tools/gen1.py (auto); review=proofread-sonnet; flags=none
\begin{question}{74}
(주)대한은 펌프사업부와 밸브사업부를 이익중심점으로 운영하고 있다. 밸브사업부는 X제품을 생산하며, X제품의 단위당 판매가격과 단위당 변동원가는 각각 ￦100과 ￦40이고, 단위당 고정원가는 ￦20이다. 펌프사업부는 연초에 Y제품을 개발했으며, Y제품을 생산하는데 필요한 A부품은 외부업체로부터 단위당 ￦70에 구입할 수 있다. 펌프사업부는 A부품 500단위를 밸브사업부로부터 대체받는 것을 고려하고 있다. 밸브사업부가 A부품 500단위를 생산 및 대체하기 위해서는 단위당 변동제조원가 ￦30과 단위당 운송비 ￦7이 발생하며, 기존 시장에서 X제품의 판매량을 200단위만큼 감소시켜야 한다. 밸브사업부가 대체거래를 수락할 수 있는 A부품의 단위당 최소 대체가격은?
\begin{choices}
  \item ￦53
  \item ￦58
  \item ￦61
  \item ￦65
  \item ￦70
\end{choices}
\end{question}
\answer{3}
